% !TeX root = ../main.tex
\section{Methodology}
\label{sec:methods}

\subsection{Dataset and preprocessing}

We use CESNET-QUIC22-XS~\cite{cesnet_quic22}, a public QUIC traffic corpus collected
on a national research and education backbone. Flows are labeled by application class.
We train on calendar week W-2022-44 and test on the subsequent week W-2022-45
(strict temporal holdout). Rare classes with fewer than 200 training samples are
removed, yielding 64 classes. Each flow is represented as a fixed-length tensor of
shape $(3, 30)$ (Fig.~\ref{fig:tensor}): channel~0 is z-scored IPT (index~0 never
jittered, following CESNET convention), channel~1 is direction, channel~2 is packet
size normalized by MTU (1500\,B). The test set contains 49{,}305 flows.

\input{figures/tikz/tensor}

\subsection{End-to-end pipeline}

Fig.~\ref{fig:pipeline} summarizes the experimental workflow. Phase~1 exports
normalized tensors; Phase~2 trains attack models; Phase~3 applies obfuscators and
logs per-flow overhead; Phase~4 evaluates frozen checkpoints and computes Pareto and
statistical summaries.

\input{figures/tikz/pipeline}

\subsection{Attack models}

We train two architectures to convergence on identical splits:
\begin{itemize}
  \item \textbf{CNN-BiLSTM:} 1D convolutions over the packet sequence followed by
        bidirectional LSTM pooling; selected hyperparameters from a six-run sweep
        (learning rate $10^{-3}$, batch size 1024, weight decay 0.01).
  \item \textbf{Masked Transformer:} multi-head self-attention encoder with padding
        mask over variable valid lengths; $d_{\mathrm{model}}=160$, trained with
        masked loss (production checkpoint used for primary results).
\end{itemize}
Checkpoints are frozen before any defense evaluation. Primary metrics are
\emph{macro-averaged F1} (class-imbalance robust) and accuracy. Chance performance
is 1.56\% (uniform over 64 classes).

\subsection{Obfuscation policies}

Obfuscators apply deterministic transforms with logged per-flow overhead:
\begin{itemize}
  \item \textbf{Jitter} (low / medium / high): add Laplace-distributed delay to
        packets with index $>0$; zero bandwidth cost, mean latency overhead 11\,ms,
        55\,ms, and 221\,ms respectively on the test manifest.
  \item \textbf{Linear-128:} round each packet size up to the next 128\,B block
        ($\approx$17.4\% mean bandwidth overhead).
  \item \textbf{Linear-128 + jitter\_medium:} compose padding and moderate jitter.
  \item \textbf{MTU:} pad every packet to 1500\,B ($\approx$274\% bandwidth overhead);
        MTU + jitter\_medium combines both mechanisms.
\end{itemize}
Seven obfuscated test tensors plus a clean baseline are evaluated by reloading weights
without retraining---modeling a defender who deploys obfuscation against a fixed,
already-trained attacker. Fig.~\ref{fig:obfuscation} illustrates the three mechanism
families.

\input{figures/tikz/obfuscation}

\subsection{Metrics and statistics}

\textbf{Privacy} is reported as classifier macro F1 and accuracy on obfuscated test
data; \textbf{cost} as mean bandwidth overhead (\%) and mean added latency (ms) from
the obfuscation manifest.

\textbf{Pareto analysis:} a setting $A$ dominates $B$ if $A$ is no worse on both cost
and privacy loss and strictly better on at least one axis. We report frontiers for
bandwidth--accuracy and latency--accuracy views.

\textbf{Tier~C statistics:} non-parametric bootstrap ($n=2000$) yields 95\% CIs for
F1 and accuracy~\cite{efron1993}. Paired bootstrap and McNemar's test~\cite{mcnemar1947}
compare baseline vs.\ \emph{jitter\_low} on matched flows. Channel ablation masks
input channels at inference without retraining.

All experiments use the open-source scripts in Phases~1--4 of the accompanying
repository; Appendix~\ref{app:repro} lists checkpoints and commands.

%==============================================================================
